\subsection{Odd gradient-square closure with trivial critical set}
\label{sec:odd-isolated-proof}

\begin{proof}[Proof of \cref{thm:odd-isolated-closure}]
The count in \eqref{eq:oddtable} applies without a nonzero critical
point. We retain its constants
\(\alpha,\beta,\gamma,\delta,\varepsilon\), without imposing the
normalization \eqref{eq:normalizedodd}. Write
\(f=F|_{\Sph^{N-1}}\), \(s=S|_{\Sph^{N-1}}\), and put
\[
 K=2e(e-1),\qquad L=4(e-1)^2,\qquad
 D=2(e-1)(N+2e-4).
\]
The spherical identities are
\begin{equation}\label{eq:odd-isolated-spherical}
\begin{aligned}
 |\nabla f|^2&=s-e^2f^2,&
 \langle\nabla f,\nabla s\rangle&=f(\beta-Ks),\\
 |\nabla s|^2&=Q(s)+\delta f^2,&
 Q(t)&=\gamma+\varepsilon t-Lt^2,\\
 \Delta f&=-e(N+e-2)f,&
 \Delta s&=\alpha-Ds.
\end{aligned}
\end{equation}
Here and throughout this proof, gradients, Hessians, and Laplacians
without subscripts are spherical. These formulas follow by subtracting
the radial derivatives from \eqref{eq:oddtable}; the Laplacian formula
uses \(\Delta(r^jv)=r^{j-2}[\Delta_{\Sph}v+j(j+N-2)v]\).

\emph{Step 1: the upper bound for \(s\).}
The assumption on the critical set gives \(s>0\). At a critical
point of \(f\), Euler's identity gives
\(\nabla_{\R^N}F=efx\), and differentiation of Euler's identity
gives \(\nabla_{\R^N}S=2(e-1)Sx\). Since \(s=e^2f^2>0\)
there, the identity for \(\Gamma(F,S)\) yields \(\beta=2e^3(e-1)f^2>0\).
Rescale \(F\) so that \(\beta=1\), and write
\[
 s_0=\frac1K,\qquad a^2=\frac{s_0}{e^2}.
\]
Thus the only critical values of \(f\) are \(a\) and \(-a\),
and \(s=s_0\) at every such point. Oddness ensures that both values
occur and that \(-a\leq f\leq a\).

We claim that \(s\leq s_0\) everywhere. Otherwise follow the complete
gradient flow of \(f\) from a point where \(g=s-s_0>0\). Along
this trajectory, with parameter \(t\),
\[
 \dot f=s-e^2f^2,\qquad \dot g=-Kfg.
\]
The second equation preserves \(g>0\). The function \(f\) increases
and is bounded, and every accumulation point is critical: indeed,
\(\int_0^\infty|\nabla f|^2\,dt<\infty\), while
\(|\nabla f|^2\) is uniformly continuous along a smooth flow on the
compact sphere. Since the initial value is larger than \(-a\), it
follows that \(f\to a\) and \(g\to0\). Direct differentiation
of the two flow equations shows that
\[
 \bigl(f^2-a^2-g/e\bigr)g^{-e/(e-1)}
\]
is constant. Consequently
\[
 f^2-a^2=\frac ge+c_1g^{e/(e-1)}.
\]
The right-hand side is positive for sufficiently small positive \(g\),
contradicting \(|f|\leq a\). This proves the claim.

\emph{Step 2: the minimum and maximum Hessians.}
Suppose that \(s\) is nonconstant. Its minimum \(m\) satisfies
\(0<m<s_0\). At a minimum point, the second identity in
\eqref{eq:odd-isolated-spherical} gives \(f=0\), so
\(v=\nabla f\) has \(|v|^2=m\). Let \(B_m=\Hess s\) there.
Differentiating that identity once, and the identity for
\(|\nabla s|^2\) twice, gives
\begin{equation}\label{eq:odd-isolated-minimum}
 B_mv=\eta v,\qquad
 B_m^2=\frac{Q'(m)}2B_m+\delta v\otimes v,
 \qquad \eta=1-Km>0.
\end{equation}
In particular,
\begin{equation}\label{eq:odd-isolated-minimum-scalar}
 Q(m)=0,\qquad
 \eta^2=\frac{Q'(m)}2\eta+\delta m.
\end{equation}
For example, the second matrix identity follows because the Hessian
of \(|\nabla s|^2\) at a critical point of \(s\) is \(2B_m^2\),
whereas the Hessian of \(\delta f^2\) at \(f=0\) is
\(2\delta v\otimes v\).

At a point where \(f=a\), put \(B=\Hess f\) and
\(C=\Hess s\). Differentiating
\eqref{eq:odd-isolated-spherical} twice gives
\begin{equation}\label{eq:odd-isolated-maximum}
 C=2B(B+e^2aI),\qquad
 BC=-e(e-1)aC,\qquad
 C^2=\left(\frac\varepsilon2-Ls_0\right)C+\delta aB.
\end{equation}
In the middle identity we used the first one to commute \(B\) and
\(C\). It follows that the eigenvalues of \(B\) belong to
\[
 0,\quad -e^2a,\quad -e(e-1)a,
\]
and the corresponding eigenvalues of \(C\) are
\(0,0,-1/e\).

\emph{Step 3: exclusion of a nonconstant \(s\).}
If \(\delta=0\), the equations at the two extrema give
\(Q(m)=Q(s_0)=0\). Hence \(Q(t)=-L(t-m)(t-s_0)\).
The scalar identity in \eqref{eq:odd-isolated-minimum-scalar},
together with \(\eta>0\), gives
\[
 K(s_0-m)=\eta=\frac{Q'(m)}2=\frac L2(s_0-m).
\]
This is impossible because \(K=2e(e-1)\ne2(e-1)^2=L/2\).

Now suppose \(\delta\ne0\). In
\eqref{eq:odd-isolated-maximum}, the eigenvalue \(-e^2a\) of
\(B\) is impossible, because its \(C\)-eigenvalue is zero.
Since \(\tr B=-e(N+e-2)a<0\), the eigenvalue
\(-e(e-1)a\) occurs. Substitution in the last identity yields
\begin{equation}\label{eq:odd-isolated-delta}
 \delta=4e-6-e\varepsilon.
\end{equation}
The value of the gradient-square identity at this same point is
\begin{equation}\label{eq:odd-isolated-q-at-maximum}
 Q(s_0)+\delta s_0/e^2=0.
\end{equation}
Set \(x=m/s_0\in(0,1)\). Eliminating \(\gamma\) with
\eqref{eq:odd-isolated-q-at-maximum}, and substituting
\eqref{eq:odd-isolated-delta}, the two scalar equations
\eqref{eq:odd-isolated-minimum-scalar}, after multiplication by
their nonzero denominators, become
\begin{align*}
 0={}&e\varepsilon[ex-(e-1)]
       +2(e^2-3e+3)-2e(e-1)x^2,\\
 0={}&e\varepsilon[ex-(e-1)]
       +2e(e-1)+(10-8e)x-2(e-1)(e-2)x^2.
\end{align*}
Their difference is
\[
 2(x-1)\bigl(2(e-1)x-(2e-3)\bigr)=0.
\]
Since \(x<1\), this gives
\begin{equation}\label{eq:odd-isolated-forced-coefficients}
 \begin{aligned}
 x&=\frac{2e-3}{2(e-1)},&
 \varepsilon&=\frac{4e^2-15e+12}{e(e-2)},\\
 \delta&=\frac e{e-2},&
 Q'(m)&=-\frac1{e-2},&
 \eta&=\frac1{2(e-1)}.
 \end{aligned}
\end{equation}
The nonnegative Hessian \(B_m\) preserves \(v^\perp\). By
\eqref{eq:odd-isolated-minimum}, its eigenvalues on \(v^\perp\)
are either zero or \(Q'(m)/2<0\), so they are all zero. Thus
\(\Delta s=\eta\) at the minimum.

At the maximum point of \(f\), let \(r\) be the multiplicity
of \(-e(e-1)a\) in \(B\). Taking traces gives
\[
 r=\frac{N+e-2}{e-1},\qquad \Delta s=-\frac re.
\]
On the other hand, subtracting the two values of
\(\Delta s=\alpha-Ds\) gives
\[
 \frac1{2(e-1)}+\frac re
 =D(s_0-m)=\frac{N+2e-4}{2e(e-1)}.
\]
Consequently \(r=(N+e-4)/(2(e-1))\), in contradiction with
the preceding trace formula. Hence \(s\) is constant, proving
\(S=cR^{e-1}\) with \(c>0\).

\emph{Step 4: the harmonic eikonal classification.}
Rescale once more so that \(S=e^2R^{e-1}\). Tkachev's polynomial
eikonal theorem \cite[Theorem~1]{TkachevIsoparametric14} gives, up to
sign,
\[
 f=T_j(g),\qquad e=jp,\qquad p\in\{1,2,3,4,6\},
\]
where \(g\) is the spherical restriction of a Cartan--M\"unzner
polynomial of degree \(p\). Since \(e\) is odd, \(p\in\{1,3\}\)
and that polynomial is harmonic. In particular,
\[
 |\nabla g|^2=p^2(1-g^2),\qquad
 \Delta g=-p(p+N-2)g.
\]
The Chebyshev equation
\((1-t^2)T_j''-tT_j'+j^2T_j=0\) and the chain rule give
\[
 \Delta f=-e^2T_j(g)-p(N-2)gT_j'(g).
\]
Comparing this with harmonicity of \(F\), namely
\(\Delta f=-e(e+N-2)f\), and using \(N>2\), yields
\[
 tT_j'(t)=jT_j(t)\qquad(-1<t<1).
\]
This identity would make \(T_j\) a monomial. For \(j>1\), its
coefficient of \(t^{j-2}\) is nonzero, so \(j=1\). Thus \(e=p=3\).
Cartan's classification of harmonic cubic eikonal polynomials
\cite{Cartan39}, also recovered in
\cite[Theorem~1.1]{TkachevCartan10}, gives the four stated
dimensions and models.
\end{proof}
